% paper3_petz_interpolation.tex
% Supplementary Material: The MOND Interpolating Function from the Petz Recovery Bound
% Supplement to Paper 3
% Author: Sheng-Kai Huang
% Compile: pdflatex paper3_petz_interpolation.tex

\documentclass[prd,twocolumn,superscriptaddress,longbibliography,floatfix]{revtex4-2}

\usepackage{amsmath}
\usepackage{amssymb}
\usepackage{amsfonts}
\usepackage{amsthm}
\usepackage{bm}
\usepackage{hyperref}
\usepackage{booktabs}
\usepackage{graphicx}
\usepackage{xcolor}

\newtheorem{theorem}{Theorem}
\newtheorem{lemma}[theorem]{Lemma}
\newtheorem{corollary}[theorem]{Corollary}
\newtheorem*{proposition}{Proposition}
\newtheorem*{remark}{Remark}

% Convenience macros
\newcommand{\kB}{k_{\mathrm{B}}}
\newcommand{\Tr}{\mathrm{Tr}}
\newcommand{\tauasym}{\tau}

% Epistemic tags
\newcommand{\known}{\textsc{[known]}}
\newcommand{\derived}{\textsc{[derived]}}
\newcommand{\assumed}{\textsc{[assumed]}}
\newcommand{\new}{\textsc{[new synthesis]}}

\begin{document}

\title{Supplementary Material: The MOND Interpolating Function\\
from the Petz Recovery Bound}

\author{Sheng-Kai Huang}
\email{akai@fawstudio.com}
\affiliation{Independent Researcher}

\date{\today}

\begin{abstract}
We show that the Petz recovery fidelity bound
$F \geq e^{-\Sigma/2}$ (Paper~1), combined with the observed
Baryonic Tully--Fisher Relation (BTFR), \emph{derives} the
deep-MOND asymptotic form $\Sigma_{\mathrm{gal}} \to 2\sqrt{x}$
($x = g_{\mathrm{bar}}/a_0$) without additional assumptions.
Extending this to all~$x$ via the minimal (zero-parameter)
extrapolation $\Sigma_{\mathrm{gal}} = 2\sqrt{x}$ yields the
interpolating function $\mu(x) = 1 - e^{-\sqrt{x}}$, which
coincides with the McGaugh--Lelli--Schombert empirical fit.
The transition-region shape ($0.1 \lesssim x \lesssim 10$) is a
\emph{genuine prediction}---not a circular restatement of BTFR.
We carefully separate what is derived from observation, what is
ansatz (the minimal extrapolation), and what constitutes a
falsifiable prediction, and identify the open problem of a
microscopic derivation of $\Sigma_{\mathrm{gal}}$ from first
principles.
\end{abstract}

\maketitle

%=====================================================================
\section{The Petz recovery bound}
\label{sec:petz}
%=====================================================================

We recall the central result of Paper~1~\cite{Huang2025paper1}.
For any completely positive trace-preserving (CPTP) channel
$\mathcal{E}$ with entropy production
\begin{equation}
  \Sigma \;=\; D\!\left(\rho \,\big\|\,
  \mathcal{R}_{\mathcal{E}}[\mathcal{E}(\rho)]\right) \geq 0\,,
  \label{eq:sigma-def}
\end{equation}
the Petz recovery fidelity satisfies the
Junge--Renner--Sutter--Wilde--Winter (JRSWW)
bound~\cite{JRSWW2018}
\begin{equation}
  F\!\left(\rho,\,
  \mathcal{R}_{\mathcal{E}}[\mathcal{E}(\rho)]\right)
  \;\geq\; e^{-\Sigma/2}\,.
  \label{eq:jrsww}
\end{equation}
Defining the temporal asymmetry parameter as the Petz recovery
failure,
\begin{equation}
  \tauasym \;=\; 1 - F\,,
  \label{eq:tau-def}
\end{equation}
the bound~\eqref{eq:jrsww} yields an upper envelope
\begin{equation}
  \tauasym_{\mathrm{bound}} \;=\;
  1 - e^{-\Sigma/2}\,.
  \label{eq:tau-bound}
\end{equation}
This is a universal, model-independent inequality: any physical
process with entropy production $\Sigma$ has its retrodictability
failure bounded by Eq.~\eqref{eq:tau-bound}. In Paper~1 we showed
that this single equation, evaluated in different limits, reproduces
the Crooks fluctuation theorem, Landauer's principle, and the
gravitational Tolman--Oppenheimer--Volkoff redshift.

%=====================================================================
\section{Galactic identification}
\label{sec:galactic}
%=====================================================================

At galactic scales, the gravitational field produces an entropy
associated with the deviation of the observed acceleration from
the Newtonian prediction. We define the standard MOND ratio
\begin{equation}
  x \;\equiv\; \frac{g_{\mathrm{bar}}}{a_0}\,,
  \label{eq:x-def}
\end{equation}
where $g_{\mathrm{bar}} = GM_{\mathrm{bar}}/r^2$ is the
Newtonian baryonic acceleration and
$a_0 \approx 1.2 \times 10^{-10}\;\mathrm{m/s}^2$ is the
MOND acceleration scale.

We identify the galactic gravitational entropy production as
\begin{equation}
  \boxed{\;\Sigma_{\mathrm{gal}} \;=\; 2\sqrt{x}
  \;=\; 2\sqrt{\frac{g_{\mathrm{bar}}}{a_0}}\;\;}
  \label{eq:sigma-gal}
\end{equation}

\textbf{Logical status of this identification.}
The deep-MOND asymptotic behavior $\Sigma_{\mathrm{gal}} \to 2\sqrt{x}$
as $x \to 0$ is \emph{derived} from the observed BTFR combined with
the Petz bound structure. The extension $\Sigma_{\mathrm{gal}} = 2\sqrt{x}$
to all $x$ is the minimal (zero-parameter) extrapolation.
The transition-region behavior of $\mu(x)$ is therefore a genuine
prediction, not a circular restatement of BTFR.
See Sec.~\ref{sec:conjecture} for a precise separation of what is
derived, what is an ansatz, and what constitutes a new prediction.

The physical motivation is as follows:
\begin{itemize}
  \item In the \emph{deep-MOND} regime ($x \to 0$): the
    gravitational field deviates maximally from Newton.
    $\Sigma_{\mathrm{gal}} \to 0$, reflecting that the
    Petz map can almost fully reconstruct the Newtonian
    state---the ``retrodiction cost'' per unit acceleration
    is small when there is almost no baryonic signal.
  \item In the \emph{Newtonian} regime ($x \to \infty$):
    $\Sigma_{\mathrm{gal}} \to \infty$, so
    $\tauasym \to 1$ (saturation), meaning that the
    Newtonian dynamics is fully ``irreversible'' at the
    information-theoretic level---no MOND correction is needed.
\end{itemize}

The $\sqrt{x}$ form is the unique power law $x^\alpha$ that
simultaneously gives $\mu \to 1$ as $x \to \infty$ (automatic
for any $\alpha > 0$) and $\mu \to \sqrt{x}$ as $x \to 0$
(requires $\alpha = 1/2$; see Sec.~\ref{sec:asymptotics}).

%=====================================================================
\section{Derivation of the interpolating function}
\label{sec:derivation}
%=====================================================================

The MOND interpolating function $\mu(x)$ relates the observed
gravitational acceleration $g_{\mathrm{obs}}$ to the baryonic
Newtonian prediction:
\begin{equation}
  g_{\mathrm{obs}} \;=\; \frac{g_{\mathrm{bar}}}{\mu(x)}\,.
  \label{eq:mond-def}
\end{equation}
In the $\tauasym$ framework, the fractional deviation of gravity
from Newton \emph{is} the Petz recovery failure: the inability to
reconstruct the Newtonian gravitational state from the observed
state quantifies precisely how much ``extra'' gravity is present.
We therefore identify
\begin{equation}
  \mu(x) \;=\; \tauasym_{\mathrm{bound}}
  \;=\; 1 - e^{-\Sigma_{\mathrm{gal}}/2}\,.
  \label{eq:mu-tau}
\end{equation}

Substituting the galactic identification~\eqref{eq:sigma-gal}:
\begin{equation}
  \boxed{\;\mu(x) \;=\; 1 - e^{-\sqrt{x}}\;\;}
  \label{eq:mu-petz}
\end{equation}

This is the main result. The MOND interpolating function arises as
the \emph{saturated Petz recovery bound}---the boundary case where
the gravitational channel is maximally efficient at converting
entropy production into retrodiction failure. Note that:
\begin{enumerate}
  \item No Boltzmann weights or thermal partition functions appear;
    this is a purely information-theoretic derivation.
  \item The only inputs are the JRSWW bound~\eqref{eq:jrsww} from
    quantum information theory and the
    identification~\eqref{eq:sigma-gal} of the galactic entropy
    production.
  \item The Crooks fluctuation theorem is \emph{not} invoked. The
    derivation lives entirely within the Petz recovery framework.
\end{enumerate}

%=====================================================================
\section{Asymptotic verification}
\label{sec:asymptotics}
%=====================================================================

\textbf{Newtonian limit} ($x \to \infty$).---As
$x \to \infty$, $e^{-\sqrt{x}} \to 0$, hence
\begin{equation}
  \mu(x) \;\to\; 1\,,
  \label{eq:newton-limit}
\end{equation}
recovering $g_{\mathrm{obs}} = g_{\mathrm{bar}}$: standard
Newtonian gravity. \checkmark

\textbf{Deep-MOND limit} ($x \to 0$).---Expanding to leading
order,
\begin{equation}
  \mu(x) = 1 - e^{-\sqrt{x}}
  = \sqrt{x} - \frac{x}{2} + \frac{x^{3/2}}{6} - \cdots
  \;\approx\; \sqrt{x}\,.
  \label{eq:mond-limit}
\end{equation}
Substituting into Eq.~\eqref{eq:mond-def}:
\begin{equation}
  g_{\mathrm{obs}} \;=\;
  \frac{g_{\mathrm{bar}}}{\sqrt{g_{\mathrm{bar}}/a_0}}
  \;=\; \sqrt{g_{\mathrm{bar}}\,a_0}\,.
  \label{eq:deep-mond}
\end{equation}
For a circular orbit, $g_{\mathrm{obs}} = V^2/r$ and
$g_{\mathrm{bar}} = GM/r^2$, giving
\begin{equation}
  V^4 \;=\; G M\,a_0\,,
  \label{eq:btfr}
\end{equation}
the Baryonic Tully--Fisher Relation (BTFR) with zero intrinsic
scatter. \checkmark

\textbf{Transition region} ($x \sim 1$).---At
$x = 1$: $\mu(1) = 1 - e^{-1} \approx 0.632$. The transition
from deep-MOND to Newtonian occurs smoothly over approximately one
decade in $x$, centered near $x \sim 1$.

%=====================================================================
\section{Comparison with other interpolating functions}
\label{sec:comparison}
%=====================================================================

Table~\ref{tab:comparison} compares the Petz-derived interpolating
function with the standard choices in the MOND literature.

\begin{table}[t]
\centering
\caption{Comparison of MOND interpolating functions.
$x = g_{\mathrm{bar}}/a_0$.}
\label{tab:comparison}
\begin{ruledtabular}
\begin{tabular}{lccc}
\textrm{Name} & $\mu(x)$ & $\mu(1)$ &
  \textrm{$x\to 0$} \\
\colrule
Petz (this work) &
  $1 - e^{-\sqrt{x}}$ & $0.632$ & $\sqrt{x}$ \\
Simple~\cite{Famaey2005} &
  $\dfrac{x}{1+x}$ & $0.500$ & $x$ \\
Standard~\cite{Milgrom1983} &
  $\dfrac{x}{\sqrt{1+x^2}}$ & $0.707$ & $x$ \\
McGaugh RAR~\cite{McGaugh2016} &
  $1 - e^{-\sqrt{x}}$ & $0.632$ & $\sqrt{x}$ \\
\end{tabular}
\end{ruledtabular}
\end{table}

Several observations:
\begin{enumerate}
  \item The Petz function~\eqref{eq:mu-petz} is
    \emph{identical in functional form} to the empirical fit
    obtained by McGaugh, Lelli, and Schombert~\cite{McGaugh2016}
    from the SPARC Radial Acceleration Relation (RAR) data.
    This is a nontrivial consistency check: the
    information-theoretic derivation independently produces the
    same function that best fits 2693 data points from 153
    galaxies.

  \item The ``simple'' and ``standard'' interpolating functions
    have $\mu \sim x$ as $x \to 0$, which gives
    $g_{\mathrm{obs}} \sim g_{\mathrm{bar}}/x
    = a_0$---a constant acceleration floor. By contrast, the
    Petz function has $\mu \sim \sqrt{x}$, giving the physically
    correct $g_{\mathrm{obs}} = \sqrt{g_{\mathrm{bar}} a_0}$.

  \item The transition sharpness differs: the Petz/McGaugh
    function transitions more gradually than the ``standard''
    function ($\mu(1) = 0.632$ vs.\ $0.707$) but more sharply
    than the ``simple'' function ($0.632$ vs.\ $0.500$).

  \item Since the Petz function matches the McGaugh RAR fit
    exactly, no additional fitting is needed against SPARC data.
    The agreement with observations is inherited directly.
\end{enumerate}

%=====================================================================
\section{Uniqueness of the Petz interpolating function}
\label{sec:uniqueness}
%=====================================================================

We now prove that the interpolating function
$\mu(x) = 1 - e^{-\sqrt{x}}$ is singled out by a combination of
asymptotic, regularity, and minimality conditions.  The results are
organized in three stages of increasing strength.

%---------------------------------------------------------------------
\subsection{Power-law uniqueness}
\label{sec:powerlaw}
%---------------------------------------------------------------------

\begin{theorem}[Power-law uniqueness]
\label{thm:powerlaw}
Let $f\colon [0,\infty) \to [0,\infty)$ be a smooth function with
$f(0) = 0$, and define $\mu(x) = 1 - e^{-f(x)/2}$.
Suppose:
\begin{itemize}
  \item[\textup{(i)}] $\mu(x) \to \sqrt{x}$ as $x \to 0^{+}$
    \textup{(deep-MOND / BTFR)};
  \item[\textup{(ii)}] $\mu(x) \to 1$ as $x \to \infty$
    \textup{(Newtonian limit)};
  \item[\textup{(iii)}] $f(x) = c\,x^{\alpha}$ for constants $c > 0$,
    $\alpha > 0$ \textup{(power-law form)}.
\end{itemize}
Then $\alpha = 1/2$ and $c = 2$.  That is,
$f(x) = 2\sqrt{x}$ and $\mu(x) = 1 - e^{-\sqrt{x}}$ are the
\emph{unique} power-law solutions.
\end{theorem}

\begin{proof}
\textbf{Step 1 (deep-MOND constraint).}
For small $x > 0$, condition~(iii) gives
$f(x)/2 = (c/2)\,x^{\alpha} \to 0$.  The Taylor expansion
$1 - e^{-u} = u - u^2/2 + \cdots$ yields
\begin{equation}
  \mu(x) \;=\; \frac{c}{2}\,x^{\alpha}
    - \frac{c^2}{8}\,x^{2\alpha} + O(x^{3\alpha})\,.
  \label{eq:mu-expand}
\end{equation}
Condition~(i) requires $\mu(x)/\sqrt{x} \to 1$ as $x \to 0^{+}$.
Comparing leading powers: $x^{\alpha} \sim x^{1/2}$, so
\begin{equation}
  \alpha = \frac{1}{2}\,.
  \label{eq:alpha-half}
\end{equation}
Matching the coefficient: $(c/2)\,x^{1/2} \to \sqrt{x}$ requires
$c/2 = 1$, hence
\begin{equation}
  c = 2\,.
  \label{eq:c-two}
\end{equation}

\textbf{Step 2 (Newtonian limit check).}
With $f(x) = 2\sqrt{x}$: as $x \to \infty$,
$e^{-\sqrt{x}} \to 0$, so $\mu(x) \to 1$.
Condition~(ii) is satisfied.~$\checkmark$

\textbf{Step 3 (uniqueness).}
Within the power-law family $f(x) = c\,x^{\alpha}$, the two
constraints~(i) and~(ii) fix both parameters uniquely.
No other member of this family satisfies both conditions.
\end{proof}

\begin{remark}
If one relaxes to $f(x) = c\,x^{\alpha} + d\,x^{\beta}$ with
$0 < \alpha < \beta$, condition~(i) still forces $\alpha = 1/2$ and
$c = 2$ (the subleading term $d\,x^{\beta}$ is irrelevant as
$x \to 0^{+}$).
Thus the deep-MOND constraint pins the leading monomial uniquely;
subleading corrections are unconstrained by asymptotics alone.
\end{remark}

%---------------------------------------------------------------------
\subsection{Strict concavity of the Petz interpolating function}
\label{sec:concavity}
%---------------------------------------------------------------------

The monotonicity and concavity of $\mu(x)$ are physically meaningful:
monotonicity says that stronger baryonic acceleration produces
dynamics closer to Newton; concavity says that the marginal return of
increasing $g_{\mathrm{bar}}$ diminishes---a kind of
``information-theoretic saturation.''

\begin{theorem}[Strict concavity]
\label{thm:concavity}
The function $\mu(x) = 1 - e^{-\sqrt{x}}$ satisfies:
\begin{itemize}
  \item[\textup{(a)}] $\mu'(x) > 0$ for all $x > 0$
    \textup{(strictly increasing)};
  \item[\textup{(b)}] $\mu''(x) < 0$ for all $x > 0$
    \textup{(strictly concave)}.
\end{itemize}
\end{theorem}

\begin{proof}
We write $\mu(x) = 1 - e^{-x^{1/2}}$ and compute all derivatives
explicitly.

\textbf{First derivative.}
Let $u(x) = x^{1/2}$, so $u'(x) = \tfrac{1}{2}\,x^{-1/2}$.
By the chain rule,
\begin{equation}
  \mu'(x) \;=\; e^{-u}\,u'
  \;=\; \frac{e^{-\sqrt{x}}}{2\sqrt{x}}\,.
  \label{eq:mu-prime}
\end{equation}
Since $e^{-\sqrt{x}} > 0$ and $\sqrt{x} > 0$ for all $x > 0$,
we have $\mu'(x) > 0$.~$\checkmark$

\textbf{Second derivative.}
We differentiate $\mu'(x) = \tfrac{1}{2}\,x^{-1/2}\,e^{-x^{1/2}}$
using the product rule.  Write
\begin{equation}
  \mu'(x) = p(x)\,q(x)\,,\quad
  p(x) = \tfrac{1}{2}\,x^{-1/2},\quad
  q(x) = e^{-x^{1/2}}\,.
\end{equation}
Their individual derivatives are
\begin{align}
  p'(x) &= -\frac{1}{4}\,x^{-3/2}\,,
  \label{eq:p-prime}\\[4pt]
  q'(x) &= e^{-x^{1/2}} \cdot \bigl(-\tfrac{1}{2}\,x^{-1/2}\bigr)
         = -\frac{e^{-\sqrt{x}}}{2\sqrt{x}}\,.
  \label{eq:q-prime}
\end{align}
By the product rule, $\mu'' = p'\,q + p\,q'$:
\begin{align}
  \mu''(x) &= \left(-\frac{1}{4x^{3/2}}\right) e^{-\sqrt{x}}
    + \frac{1}{2\sqrt{x}} \cdot
      \left(-\frac{e^{-\sqrt{x}}}{2\sqrt{x}}\right)
    \nonumber\\[4pt]
  &= -\frac{e^{-\sqrt{x}}}{4\,x^{3/2}}
     -\frac{e^{-\sqrt{x}}}{4\,x}\,.
  \label{eq:mu-double-prime-expand}
\end{align}
Factoring out $-e^{-\sqrt{x}}/4$:
\begin{equation}
  \boxed{\;\mu''(x) \;=\;
    -\frac{e^{-\sqrt{x}}}{4}
    \left(\frac{1}{x^{3/2}} + \frac{1}{x}\right)\;\;}
  \label{eq:mu-double-prime}
\end{equation}
For $x > 0$: $e^{-\sqrt{x}} > 0$, $x^{-3/2} > 0$, $x^{-1} > 0$,
so the factor in parentheses is strictly positive, and the overall
sign is negative.  Therefore $\mu''(x) < 0$ for all $x > 0$.
\end{proof}

\begin{remark}
The two terms in~\eqref{eq:mu-double-prime} have distinct origins.
The $x^{-3/2}$ term comes from the curvature of the square-root
mapping $x \mapsto \sqrt{x}$ (i.e.\ from $p'$), while the $x^{-1}$
term comes from the exponential saturation of the Petz bound (i.e.\
from $q'$).
Both contribute with the same sign, so concavity is \emph{robust}:
it does not rely on a cancellation.
\end{remark}

%---------------------------------------------------------------------
\subsection{Characterization theorem}
\label{sec:characterization}
%---------------------------------------------------------------------

The following theorem combines the results above into a single
characterization.

\begin{theorem}[Characterization]
\label{thm:characterization}
Let $f\colon [0,\infty) \to [0,\infty)$ be smooth and monotonically
increasing with $f(0) = 0$, and define
$\mu_{f}(x) = 1 - e^{-f(x)/2}$.  Suppose:
\begin{itemize}
  \item[\textup{(i)}] $\mu_{f}(x)/\sqrt{x} \to 1$ as $x \to 0^{+}$
    \textup{(BTFR)};
  \item[\textup{(ii)}] $\mu_{f}(x) \to 1$ as $x \to \infty$
    \textup{(Newtonian limit)}.
\end{itemize}
Then:
\begin{enumerate}
  \item \textbf{Asymptotic form.}
    $f(x) = 2\sqrt{x} + h(x)$, where $h(x)/\sqrt{x} \to 0$
    as $x \to 0^{+}$.
  \item \textbf{Monomial uniqueness.}
    If $f$ is a monomial ($f(x) = c\,x^{\alpha}$), then
    $f(x) = 2\sqrt{x}$.
  \item \textbf{Concavity of the minimal solution.}
    The minimal (zero-parameter) choice $h = 0$, giving
    $\mu(x) = 1 - e^{-\sqrt{x}}$, is strictly concave on
    $(0,\infty)$.
  \item \textbf{Zero-parameter uniqueness.}
    Among all functions of the form
    $\mu_{f}(x) = 1 - e^{-f(x)/2}$ satisfying \textup{(i)}
    and~\textup{(ii)}, the choice $f(x) = 2\sqrt{x}$ is the unique
    member with no free parameters beyond those fixed by
    \textup{(i)} and~\textup{(ii)}.
\end{enumerate}
\end{theorem}

\begin{proof}
\textbf{Part 1.}
Condition~(i) with $f(x) \to 0$ gives
$\mu_{f}(x) \approx f(x)/2 + O(f^2)$.  Hence $f(x)/2 \to \sqrt{x}$
as $x \to 0^{+}$, i.e.\ $f(x) = 2\sqrt{x} + h(x)$ with
$h(x)/\sqrt{x} \to 0$.

\textbf{Part 2.}
If $f(x) = c\,x^{\alpha}$, Part~1 forces the leading behavior
$c\,x^{\alpha} \sim 2\,x^{1/2}$, so $\alpha = 1/2$ and $c = 2$.
This is Theorem~\ref{thm:powerlaw}.

\textbf{Part 3.}
This is Theorem~\ref{thm:concavity}.

\textbf{Part 4.}
Any $f$ satisfying Part~1 takes the form $f(x) = 2\sqrt{x} + h(x)$
with $h(x)/\sqrt{x} \to 0$.  The correction~$h$ is a free function
(subject to regularity).  Setting $h = 0$ is the unique choice that
introduces no additional functional freedom.  Every other choice
requires specifying the form of~$h$ (e.g.\ $h(x) = d\,x^{\beta}$
with $\beta > 1/2$), which introduces at least one free parameter.
\end{proof}

\begin{corollary}
\label{cor:occam}
If one imposes Occam's razor---the interpolating function
$\mu(x) = 1 - e^{-f(x)/2}$ should involve no free parameters beyond
those fixed by the BTFR and Newtonian asymptotics---then
$\mu(x) = 1 - e^{-\sqrt{x}}$ is the unique solution.  This
zero-parameter function is automatically strictly monotone and
strictly concave.
\end{corollary}

%=====================================================================
\section{What remains conjectural}
\label{sec:conjecture}
%=====================================================================

It is important to clearly separate what is derived from observation,
what is an extrapolation, and what constitutes a genuinely new prediction.
We organize the logical status of each element below.

\subsection{Derived from BTFR (not conjectural)}

The deep-MOND asymptotic form $\Sigma_{\mathrm{gal}} \to 2\sqrt{x}$
as $x \to 0$ is \derived, not \assumed. The argument is as follows.
Within the Petz bound framework, $\mu(x) = 1 - e^{-\Sigma_{\mathrm{gal}}/2}$.
In the deep-MOND limit ($x \to 0$), $\Sigma_{\mathrm{gal}} \to 0$
and $\mu \approx \Sigma_{\mathrm{gal}}/2$.
The BTFR ($V^4 = GMa_0$) requires $\mu \to \sqrt{x}$ as $x \to 0$
(see Eq.~\eqref{eq:btfr}).
Combining: $\Sigma_{\mathrm{gal}}/2 \to \sqrt{x}$, i.e.\
$\Sigma_{\mathrm{gal}} \to 2\sqrt{x}$.

This step uses only: (a)~the JRSWW bound structure, and
(b)~the empirically established BTFR.
It is a \emph{derivation}, not a fit.

\subsection{Ansatz: minimal zero-parameter extrapolation}

The extension
\begin{equation}
  \Sigma_{\mathrm{gal}} = 2\sqrt{x}
  \quad \text{for \emph{all} } x
  \tag{\ref{eq:sigma-gal}}
\end{equation}
is the simplest (zero-parameter) extrapolation of the derived
deep-MOND asymptotic form to the full range of $x$.
This is an \assumed\ ansatz, not a derivation.
More complex functional forms---e.g.\
$\Sigma_{\mathrm{gal}} = 2\sqrt{x}\,h(x)$ with $h(x) \to 1$ as
$x \to 0$---could in principle modify the transition region while
preserving the BTFR.
The choice $h(x) = 1$ (no additional parameters) is the minimal
hypothesis.

Three supporting arguments for this minimal choice:

\textbf{(i) Dimensional consistency.}
$\Sigma$ must be dimensionless. The combination
$g_{\mathrm{bar}}/a_0$ is dimensionless, and its square root is
the simplest function thereof. The prefactor~2 is fixed by matching
the JRSWW bound normalization~\eqref{eq:jrsww}.

\textbf{(ii) Power-law uniqueness.}
Consider the general power-law ansatz $\Sigma_{\mathrm{gal}} = 2\,x^\alpha$.
Then $\mu = 1 - e^{-x^\alpha}$. The deep-MOND limit
$\mu \approx x^\alpha$ must equal $\sqrt{x}$ to reproduce the
BTFR~\eqref{eq:btfr}, forcing $\alpha = 1/2$. No other power law
works. (Note: this fixes the \emph{exponent} but does not exclude
non-power-law corrections at finite~$x$.)

\textbf{(iii) Connection to the cosmological sector.}
If $\mu_0 = H_0/c$ (as suggested by Paper~3's cosmological
analysis) and $a_0 = cH_0/(2\pi)$ (Milgrom's coincidence), then
\begin{equation}
  \Sigma_{\mathrm{gal}} = 2\sqrt{\frac{g_{\mathrm{bar}}}
  {c^2\mu_0/(2\pi)}}\,,
  \label{eq:sigma-cosmo}
\end{equation}
linking the galactic entropy production to the same parameter
$\mu_0$ that governs the cosmological sector.

\subsection{Genuinely new prediction}

Given the minimal ansatz $\Sigma_{\mathrm{gal}} = 2\sqrt{x}$ for
all~$x$, the interpolating function
$\mu(x) = 1 - e^{-\sqrt{x}}$ makes a \emph{genuine prediction} for
the transition-region behavior at $0.1 \lesssim x \lesssim 10$,
where neither the deep-MOND nor the Newtonian asymptotic applies.
This prediction is \emph{not} a circular restatement of BTFR (which
constrains only $x \to 0$) nor of Newtonian gravity (which constrains
only $x \to \infty$).
The transition shape---in particular, $\mu(1) = 1 - e^{-1} \approx 0.632$---is
testable against SPARC RAR data and constitutes the falsifiable
content of the ansatz.

The fact that this predicted shape coincides with the McGaugh et
al.~\cite{McGaugh2016} empirical fit to 2693 data points from 153
galaxies is a nontrivial success of the minimal extrapolation, not
a built-in feature.

\subsection{Open problem: microscopic derivation}

A microscopic derivation would require computing the quantum
relative entropy $D(\rho_{\mathrm{obs}}\|\rho_{\mathrm{Newton}})$
for the gravitational field in the weak-field, quasi-static limit,
and showing that $\Sigma_{\mathrm{gal}} = 2\sqrt{x}$ emerges
without the minimal-extrapolation assumption. This would promote
the ansatz to a theorem. We leave this to future work.

%=====================================================================
\section{Summary}
\label{sec:summary}
%=====================================================================

The derivation chain is:
\begin{equation}
  \underbrace{F \geq e^{-\Sigma/2}}_{\text{JRSWW bound}}
  \;\xrightarrow{\;\Sigma_{\mathrm{gal}} = 2\sqrt{x}\;}\;
  \underbrace{\mu(x) = 1 - e^{-\sqrt{x}}}_{\text{MOND interpolation}}
  \label{eq:chain}
\end{equation}
The Petz recovery bound, combined with a single identification of
the galactic entropy production, yields the MOND interpolating
function that coincides with the best empirical fit to the Radial
Acceleration Relation. The derivation uses no Boltzmann weights, no
Crooks theorem, and no partition functions---only the
information-theoretic structure of quantum channels.

\begin{acknowledgments}
This work builds on the $\tauasym$ framework established in
Paper~1~\cite{Huang2025paper1}.
\end{acknowledgments}

\begin{thebibliography}{9}

\bibitem{Huang2025paper1}
S.-K.~Huang,
``Temporal asymmetry as Petz recovery failure: a unifying
framework from the Crooks relation to gravitational redshift,''
(2025); \href{https://doi.org/10.5281/zenodo.14887770}
{Zenodo: 10.5281/zenodo.14887770}.

\bibitem{JRSWW2018}
M.~Junge, R.~Renner, D.~Sutter, M.~M.~Wilde, and A.~Winter,
``Universal recovery maps and approximate sufficiency of quantum
relative entropy,''
Ann.\ Henri Poincar\'{e} \textbf{19}, 2955 (2018).

\bibitem{Crooks1999}
G.~E.~Crooks,
``Entropy production fluctuation theorem and the
nonequilibrium work relation for free energy differences,''
Phys.\ Rev.\ E \textbf{60}, 2721 (1999).

\bibitem{McGaugh2016}
S.~S.~McGaugh, F.~Lelli, and J.~M.~Schombert,
``Radial Acceleration Relation in Rotationally Supported
Galaxies,''
Phys.\ Rev.\ Lett.\ \textbf{117}, 201101 (2016).

\bibitem{Famaey2005}
B.~Famaey and J.~Binney,
``Modified Newtonian dynamics in the Milky Way,''
Mon.\ Not.\ R.\ Astron.\ Soc.\ \textbf{363}, 603 (2005).

\bibitem{Milgrom1983}
M.~Milgrom,
``A modification of the Newtonian dynamics as a possible
alternative to the hidden mass hypothesis,''
Astrophys.\ J.\ \textbf{270}, 365 (1983).

\end{thebibliography}

\end{document}
